\documentclass[a4paper, 10pt]{style/preprint}
\usepackage[left=1.5in, right=1.5in, bottom=1.5in]{geometry}
\usepackage{style/comment}
\usepackage{style/mhenvs}
\usepackage{style/mhsymb}

\usepackage[full]{textcomp}
\usepackage[osf]{newtxtext}
\usepackage{amssymb}
\usepackage{amsmath}
\usepackage{mathtools}
\usepackage{mhequ}
\usepackage{booktabs}
\usepackage{tikz}
\usepackage{tikz-cd}
\usepackage{mathrsfs}
\usepackage{longtable}
\usepackage{microtype}
\usepackage{wasysym}
\usepackage{centernot}
\usepackage{enumitem}
\usepackage[dvipsnames]{xcolor}
\usepackage{hyperref}
\hypersetup{
  pdftitle={Statistics Notes},
  pdfauthor={Kexing Ying},
  colorlinks=true,
  linkcolor=Maroon,
  filecolor=Maroon,
  citecolor=Maroon,
  urlcolor=Maroon,
  pdfcreator={LaTeX via pandoc}}
\usepackage{crossreftools}
\usepackage{physics}
\usepackage[notref,notcite]{showkeys}

\input{style/prelude.tex}

\DeclareMathOperator{\EPF}{EPE}
\def\pvar{p\mathrm{\text{-}var}}

\begin{document}

\title{Statistics}
\author{Kexing Ying}
\institute{}
\maketitle

\section{Regression function}

For square loss function \(L(Y, f(X)) = (Y - f(X))^2\), the expected prediction error is
\begin{equ}
  \EPF(f) = \mathbb{E}[(Y - f(X))^2].
\end{equ}
It is easy to see that the regression function \(f(x) = \mathbb{E}[Y \mid X = x]\) minimizes the
expected prediction error. Indeed, as all \(\sigma(X)\)-measurable random variables are of the form
\(f(X)\), it is equivalent to minimize \(\|Y - Z\|_2\) over all \(\sigma(X)\)-measurable random
variables \(Z\). This is of course achieved by the orthogonal projection of \(Y\) onto the closed
subspace of \(\sigma(X)\)-measurable functions which is exactly \(\mathbb{E}[Y \mid X]\).

The nearest neighbour method directly tries to estimate this regression function. In particular,
intuitively, given iid \((X_i, Y_i)_{i = 1}^n\) and \(x\), we hope to approximate
\begin{equ}
  \mathbb{E}[Y \mid X = x] \approx \frac{1}{k} \sum_{i = 1}^n \mathbf{1}_{\{\|X_i - x\| \le \|X_{(k)} - x\|\}} Y_{i}.
\end{equ}
Namely, we take the sample average of the \(Y_i\) values corresponding to the \(k\) nearest neighbours of \(x\).
By taking a variance-bias decomposition, if we take \(k, n\) in such a way that \(k \to \infty\) and
\(\frac{k}{n} \to 0\), then under some mild assumptions on the distribution of \((X, Y)\), we have
the approximation converges in the limit.

The nearest neighbour method (and local methods in general) is problematic in high dimensions.
For fixed \(x\), if we want our estimate to depend on the data within a neighbourhood of \(x\) which
captures a fixed proportion of the data, then the radius of this neighbourhood has to grow
exponentially in the dimension (so rather contradicts the intuition of locality). This is called
\textit{the curse of dimensionality}.

A natural extension of the nearest neighbour method is, instead of taking the average of the nearest
neighbours, we can take a weighted average of all points weighted by their distance to \(x\).
We call the weight function \(K\) the kernel function and a simple estimate (Nadaraya--Watson) is
\begin{equ}
  f(x) = \frac{1}{Z}\sum_{i=1}^n K\left(x, X_i\right) Y_i
\end{equ}
with the normalisation \(Z = \sum_{i=1}^n K(x, X_i)\).

\section{Linear regression}

Let us consider the OLS estimator from a probabilistic perspective. Suppose \(\mathcal{H}, \mathcal{K}\)
are two Hilbert spaces and we have random variables \(X \in \mathcal{H}\) and \(Y \in \mathcal{K}\).
The fundamental assumption one makes with linear regression is that the regression function is linear,
i.e.\ there exists a linear operator \(\beta : \mathcal{H} \to \mathcal{K}\) such that
\(\mathbb{E}[Y \mid X] = \beta X\) (we assume w.l.o.g.\ that both \(X, Y\) are centred).

Now, recalling that the covariance operators \(C_X\) and \(C_{XY}\) are defined by
\begin{equ}
  C_X : \mathcal{H} \to \mathcal{H}, \quad C_X f = \mathbb{E}[\scal{f, X}_\mathcal{H} X]
\end{equ}
and
\begin{equ}
  C_{XY} : \mathcal{H} \to \mathcal{K}, \quad C_{XY} f = \mathbb{E}[\scal{f, X}_\mathcal{H} Y].
\end{equ}
Thus, we compute
\begin{equ}
  \scal{C_{XY}h, k} = \mathbf{E}[\scal{X, h}\scal{Y, k}] = \mathbf{E}[\scal{X, h}\scal{\mathbf{E}[Y \mid X], k}]
  = \mathbf{E}[\scal{X, h}\scal{\beta X, k}] = \scal{\beta C_{X} h, k}
\end{equ}
for all \(h \in \mathcal{H}\) and \(k \in \mathcal{K}\). Hence, we have \(C_{XY} = \beta C_X\).
Thus, in the case where \(C_X\) is invertible, we have \(\beta = C_{XY} C_X^{-1}\).

In the finite-dimensional case, the OLS estimator is exactly the empirical version of this formula.
In particular, suppose now \(\mathcal{H} = \mathbb{R}^p\) and \(\mathcal{K} = \mathbb{R}\).
Then, we have the covariance matrices \(C_X = \mathbb{E}[X X^\top]\) and \(C_{XY} = \mathbb{E}[X Y^\top]\).
Now, given iid data \((\mathbf{X}, \mathbf{Y}) = (X_i, Y_i)_{i = 1}^n \in \mathbb{R}^{n \times p} \times \mathbb{R}^{n}\), we can estimate \(C_X\) and \(C_{XY}\) by their empirical versions
\begin{equ}
  \hat{C}_X = \frac{1}{n} \sum_{i = 1}^n X_i X_i^\top = \frac{1}{n} \mathbf{X}^\top \mathbf{X},
  \quad \hat{C}_{XY} = \frac{1}{n} \sum_{i = 1}^n X_i Y_i^\top = \frac{1}{n} \mathbf{X}^\top \mathbf{Y}.
\end{equ}
Thus, the OLS estimator is exactly \(\hat{\beta} = \hat{C}_{XY} \hat{C}_X^{-1}\). Note that the OLS
estimator one typically writes down is the transpose of what we have here as one typically instead assumes
\(Y = X \beta + \epsilon\).

\section{Ridge regression and kernels}

The OLS estimator is the unbiased estimator minimising MSE. However, it is not the estimator minimising MSE
among all estimators. In particular, via the bias-variance decomposition, one might hope to reduce the
MSE by introducing some bias. The ridge regression estimator achieves this by penalising large values of \(\beta\).
This is intuitive as \(\beta\) is large when \(C_X\) is close to singular and thus the OLS estimator is unstable.
The ridge regression estimator instead optimises
\begin{equ}
  \hat{\beta} = \arg\min_{\beta} \left( \frac{1}{n} \sum_{i=1}^n (Y_i - X_i^\top \beta)^2 + \lambda \|\beta\|^2 \right)
\end{equ}
for some regularisation parameter \(\lambda > 0\).

\subsection{The Cameron--Martin space and RKHS}

Suppose now \(X\) is a Gaussian random variable in a Hilbert space \(\mathcal{B}\) with law \(\mu\).
We recall that the Cameron--Martin space \(\mathcal{H}_\mu\) of \(\mu\) is defined as the closure of
the pre-Cameron--Martin space
\begin{equ}
  \mathcal{H}_{\mu}^\circ = \left\{ x \in \mathcal{B} : \exists\, x^{*} \in \mathcal{B}^*,\
  l(x) = C_\mu(l, x^{*})\ \forall\, l \in \mathcal{B}^{*} \right\} = C_\mu(\mathcal{B}^*)
\end{equ}
where in the latter equality we identify \(C_\mu : \mathcal{B}^* \to \mathcal{B}\) by
\(C_\mu l = \int l(x) x\, \mu(\dd x)\) with the integral understood in the Bochner sense.

We show that the pre-Cameron--Martin space is exactly the \RKHS{} (in the statistician's sense) associated to the kernel
\(k(h, k) = C_\mu(h, k) = \int h(x) k(x)\, \mu(\dd x)\) (to be more precise, the \RKHS{} of \(k\) is
the canonical identification of \(\mathcal{H}_\mu^\circ\) in \(\mathcal{B}^{**}\)).

Recall that the \RKHS{} \(\mathcal{H}\) associated to a kernel
\(k : \mathcal{X} \times \mathcal{X} \to \mathbb{R}\) is a Hilbert space of functions
\(f : \mathcal{X} \to \mathbb{R}\) such that \(k(x, \cdot) =: k_x \in \mathcal{H}\) for all \(x\) and
\begin{equ}
  \scal{f, k(x, \cdot)}_{\mathcal{H}} = f(x).
\end{equ}
In this case, \(\mathcal{X} = \mathcal{B}^*\) and \(\mathcal{H} = \mathcal{H}_\mu\) with the canonical identification
\(x \in \mathcal{H}_\mu \subseteq \mathcal{B}^{**}\), \(x: l \mapsto l(x)\). Indeed, in this case,
\(k(l, \cdot) = C_\mu l\) and for any \(x \in \mathcal{H}_\mu\), we have that
\begin{equ}
  \scal{x, k(l, \cdot)}_\mu = \scal{x, C_\mu l}_\mu = C_\mu(x^{*}, l) = l(x) = x(l)
\end{equ}
where we used the fact that \((C_\mu l)^* = l\).

\bibliographystyle{alpha}
\bibliography{refs}

\end{document}
